¿Quiénes consumen Kali Linux y con qué finalidad?

Kali Linux es principalmente utilizada por profesionales de ciberseguridad como pentesters, auditores y consultores que se encargan de evaluar la robustez de los sistemas de las organizaciones. Además, se utiliza en equipos de seguridad (\textit{Red Teams} y \textit{Blue Teams}) que se encargan de hacer simulación de ataques para conocer los métodos de ataque y de esta manera mejorar las defensas de los sistemas.

Por las herramientas de análisis de red como Wireshark, es común que administradores de sistemas y redes lo utilicen para verificar la robustez de redes, servidores y las configuraciones de seguridad de estos entornos.

Finalmente, dado que muchas de sus herramientas también pueden ser usadas de manera maliciosa, esta distribución también es utilizada por atacantes y ciberdelincuentes para explotar vulnerabilidades o encontrar fallas en los sistemas.

En resumen, Kali Linux se utiliza para múltiples tareas relacionadas con la seguridad como la recopilación de información de sistemas, el análisis de vulnerabilidades y pruebas de penetración, así como para realizar ataques de prueba, idealmente simulados y con la debida autorización del propietario del sistema.
